\documentclass[11pt]{article}
\usepackage[utf8]{inputenc}
\usepackage{hyperref}
\usepackage{geometry}
\geometry{margin=1in}

\title{ADN-Chain v1: Codon-Based Virtual Machine and Blockchain System}
\author{Jamil Al Thani}
\date{Version 1.0.0}

\begin{document}
\maketitle
\begin{abstract}
ADN-Chain v1 is a codon-inspired distributed computing platform that combines a
deterministic virtual machine, a 64-codon instruction set, and a blockchain
data layer. This whitepaper outlines the core design and implementation
considerations for researchers and contributors.
\end{abstract}

\section{Introduction}
This document introduces the motivations behind ADN-Chain v1, the biological
inspiration for codon-based computation, and the objectives of the initial
release.

\section{Computation Model}
We describe the execution semantics, memory model, and determinism guarantees
that underpin the ADN virtual machine family.

\section{ISA-64}
The 64-codon instruction set encodes arithmetic, control flow, memory access,
and system operations. Each codon maps to a mnemonic and semantics documented in
this section.

\section{Bytecode}
The bytecode format specifies module layout, symbol tables, and relocation rules
needed to produce portable codon programs.

\section{Virtual Machine}
We outline the ADN-VM and CodonVM, including interpreter pipelines, runtime
state, host interfaces, and deterministic execution constraints.

\section{ABI and Syscalls}
This section defines the Application Binary Interface, syscall catalog, calling
conventions, and safety rules for contract and node integrations.

\section{State and Transactions}
We describe the state model, account and storage layout, transaction types, and
state transition function for block validation.

\section{Blocks}
Blocks bundle transactions, execution receipts, and cryptographic commitments.
We cover block headers, Merkle structures, and validation rules.

\section{Network}
ADN-Net v1 is a P2P protocol for propagating blocks, transactions, and state
snapshots. Messaging primitives and bandwidth considerations are summarized
here.

\section{Consensus}
This release targets a deterministic, slot-based consensus inspired by D10Z.
The section details validator roles, leader selection, and fork choice.

\section{Cryptography}
We specify key types, signature schemes, hashing primitives, and encoding
formats used across the stack.

\section{Implementation Notes}
Implementation guidance for Rust crates, workspace layout, testing strategy, and
interoperability with WASM toolchains.

\section{Toolchain}
Build tooling, CLI utilities, and developer workflows for authoring and testing
codon programs.

\section{Roadmap}
Future work spans performance optimizations, formal verification, and expanded
network protocols. A milestone roadmap is provided for contributors.

\section{Conclusion}
ADN-Chain v1 establishes a foundation for codon-based distributed computation.
Community contributions are encouraged to mature the protocol and tooling.

\end{document}
